\chapter{Градуированное соответствие, суперсвязность и калибровочная развилка}

Коммутантный гейт доказал, что обычные одноформы алгебры Стандартной модели
не видят семейный оператор. Поэтому проверим более широкий, но всё ещё
строгий объект: градуированное соответствие со связностью. В языке
неограниченной бивариантной $K$-теории соответствие с оператором и
связностью задаёт морфизм спектральных геометрий, а произведение имеет вид
\begin{equation}
 D_{\mathrm{tot}}
 =S\widehat\otimes1+1\widehat\otimes_{\nabla}D_{\mathrm{SM}}.
 \label{eq:v5-correspondence-product}
\end{equation}
Именно такой объект допускает вертикальный семейный оператор, не объявляя
семейную матричную алгебру второй координатной алгеброй.

\section{Минимальное соответствие}

Возьмём правый модуль
\begin{equation}
 E=K_{\mathrm{fam}}\widehat\otimes\mathcal A_{\mathrm{SM}},
 \qquad
 K_{\mathrm{fam}}=\mathbb C^4\oplus\mathbb C^3\oplus\mathbb C^3,
\end{equation}
и семейный нечётный эндоморфизм
\begin{equation}
 Q_{\mathrm{fam}}=d+d^\dagger.
\end{equation}
Поскольку $Q_{\mathrm{fam}}$ действует только на первом множителе,
\begin{equation}
 [Q_{\mathrm{fam}},R(a)]=0,
 \qquad a\in\mathcal A_{\mathrm{SM}}.
\end{equation}
Следовательно,
\begin{equation}
 \nabla=\nabla_{\mathrm{SM}}+Q_{\mathrm{fam}}
\end{equation}
удовлетворяет градуированному правилу Лейбница. Численный дефект
$\mathcal A_{\mathrm{SM}}$-линейности меньше $7\cdot10^{-16}$. Условия KO6
на семейной части наследуются от построенного ранее двадцатимерного
бимодуля.

Это настоящий положительный результат: семейная цепь допустима как данные
морфизма. Но она является дополнительным операторным содержанием
соответствия, а не следствием одной координатной алгебры.

\section{Три кривизны и возвращение проблемы знака}

Для цепи $4\to3\to3$ рассмотрим три естественных кандидата.

\subsection{Самосопряжённая суперсвязность}

Стандартная нулевая по форме часть кривизны равна
\begin{equation}
 Q_{\mathrm{fam}}^2=(d+d^\dagger)^2.
\end{equation}
Её средний блок есть
\begin{equation}
 Q_{\mathrm{fam}}^2\big|_{V_G}
 =XX^\dagger+Y^\dagger Y
 =XX^\dagger+|\Phi|^2I_3.
 \label{eq:v5-superconnection-gram-sum}
\end{equation}
Таким образом, обычная суперсвязность воспроизводит уже известную сумму
матриц Грама, а не отображение момента.

\subsection{Односторонняя суперсвязность}

Если оставить только повышающий оператор $d$, то
\begin{equation}
 d^2:V_L\longrightarrow V_R,
 \qquad
 d^2\big|_{V_L}=\Phi X.
\end{equation}
Его положительная норма равна
\begin{equation}
 \Tr (d^2)^\dagger d^2
 =|\Phi|^2\Tr X^\dagger X.
 \label{eq:v5-holomorphic-path-norm}
\end{equation}
Она видит составной путь, но также не даёт
$XX^\dagger-|\Phi|^2I_3$.

\subsection{Эрмитово отображение момента}

Искомая величина возникает только как
\begin{equation}
 \mathcal F_H=[d,d^\dagger]-h.
 \label{eq:v5-hamiltonian-hodge-curvature}
\end{equation}
Это каноническое эрмитово отображение момента ориентированного колчана, но
не стандартный квадрат суперсвязности. Следовательно, градуированное
соответствие само по себе не устраняет прежний разрыв. Оно должно быть
обогащено гамильтоновой, или эквивалентной ей эрмитово-моментной,
структурой.

\begin{theorem}[Предел обычной суперсвязности]
На аффинной цепи стандартная самосопряжённая суперсвязность даёт сумму
$XX^\dagger+Y^\dagger Y$, а односторонняя --- составной путь $YX$.
Разность $XX^\dagger-Y^\dagger Y$ требует отдельной эрмитовой структуры
отображения момента.
\end{theorem}

Этот вывод согласуется с литературой о неограниченных соответствиях:
связность является самостоятельной частью морфизма, а внутренние флуктуации
распадаются на связности и эндоморфизмы модуля. Литература обеспечивает
язык, но не выводит конкретное проектное отображение момента.

\section{Поправка к пространству допустимых флуктуаций}

Аффинная высота имеет вид
\begin{equation}
 h=\operatorname{diag}(-P_3,0_3,I_3).
\end{equation}
В вакууме $X=V$ выполнено
\begin{equation}
 [h,d]=d.
\end{equation}
Однако для произвольной вариации левой стрелки
\begin{equation}
 [h,\delta d]=\delta d
 \quad\Longleftrightarrow\quad
 \delta X=\delta XP_3.
 \label{eq:v5-fixed-height-tangent}
\end{equation}
Три вещественные вариации $\delta XP_1$ нарушают условие степени; их дефект
равен единице в нормированном базисе.

Это уточняет предыдущий гессиан. Три собственных значения $4/3$, найденные
для смешивания опорного синглета с цепью, действительно положительны в
расширенном функционале Ходжа, но эти направления не являются касательными
к пространству суперсвязностей с фиксированной высотой. В строгом
градуированном соответствии они отсутствуют, а не получают массу.

\section{Что исправляет общая нормировка}

Для любой вариации семейного оператора полный след даёт
\begin{equation}
 \tau_{300}(\delta Q^\dagger\delta Q)
 =\frac1{10}\Tr_{10}(\delta Q^\dagger\delta Q)
 =\frac3{10}G_{\mathrm{red}}(\delta Q,\delta Q).
\end{equation}
Тот же множитель $3/10$ умножает квадрат кривизны. Максимальные численные
дефекты этих двух тождеств равны соответственно
$1.8\cdot10^{-15}$ и $1.5\cdot10^{-14}$.

Следовательно, общий множитель сокращается в обобщённой задаче
\begin{equation}
 H v=m^2Gv.
\end{equation}
Одна только поправка нормировки полного следа не меняет вакуум, знак и
массовые собственные значения сокращённой модели. Открытым остаётся не
множитель $3/10$, а вывод относительных коэффициентов кинетического и
потенциального членов из одного полного действия.

\section{Калибровочная и БВ/БРСТ-развилка}

Если калибровочной остаётся только группа Стандартной модели, семейный
$SO(3)$ не является калибровочной симметрией. Тогда три вращательных нуля
гессиана не являются БРСТ-точными и представляют физические безмассовые
моды.

Если же калибровать все унитарные эндоморфизмы кратностного модуля, то уже
минимальный подгрупповой множитель содержит $U(10)$ размерности $100$.
Возникают новые калибровочные поля и духи, что противоречит условию
отсутствия калибровочного расширения.

Поэтому нельзя одновременно утверждать:
\begin{enumerate}
 \item что семейный $SO(3)$ не создаёт новых бозонов;
 \item что три его орбитальных нуля удаляются калибровочным
 факторированием.
\end{enumerate}

\begin{nogocriterion}[Коммутантное БВ/БРСТ-замыкание]
В ветви с одной калибровочной алгеброй Стандартной модели семейные
вращательные нули нельзя удалять как калибровочные. В ветви с калибровкой
эндоморфизмов соответствия возникает запрещённое расширение группы.
\end{nogocriterion}

\section{Возврат к старому тетраэдрическому сектору}

Ранние главы проекта уже содержат объект, способный снять эту развилку без
непрерывных семейных бозонов: тензор третьего ранга
\begin{equation}
 \mathcal T_{ijk}=\sum_{a=1}^4n_{a,i}n_{a,j}n_{a,k},
 \qquad
 \operatorname{Stab}_{SO(3)}(\mathcal T)=A_4.
\end{equation}
Он дискретно нарушает случайную симметрию $SO(3)$ до $A_4$ и уже связан в
проекте с кубическим корневым расслоением и голономией $\mathbb Z_3$.

Это не готовое спасение: пока не доказано, что тетраэдрический член входит
в тот же след $\tau_{300}$ с фиксированным коэффициентом. Но это первый
следующий шаг, который использует уже выведенный объект проекта, снимает
вращательные нули без новых непрерывных калибровочных полей и не возвращает
наивное тензорное произведение алгебр.

\begin{protocol}
\begin{enumerate}
 \item градуированное правило Лейбница:
 \textbf{выполнено};
 \item KO6 и ковариантность относительно Стандартной модели:
 \textbf{выполнены};
 \item стандартная кривизна суперсвязности как отображение момента:
 \textbf{не проходит};
 \item общий след для кинетической и потенциальной нормировок:
 \textbf{проходит на уровне общего множителя};
 \item три опорно-цепные моды внутри фиксированной градуировки:
 \textbf{отсутствуют};
 \item удаление семейных нулей при одной группе Стандартной модели:
 \textbf{не проходит};
 \item калибровка всех эндоморфизмов соответствия:
 \textbf{закрыта расширением группы};
 \item дискретное тетраэдрическое снятие нулей:
 \textbf{открыто для следующего гейта};
 \item физическое замыкание:
 \textbf{не достигнуто}.
\end{enumerate}
\end{protocol}

Машинная проверка и сертификат сохранены в
\path{s2t_v5_graded_correspondence_superconnection_gate.py} и
\path{s2t_v5_graded_correspondence_superconnection_gate_results.json}.